\documentclass[12pt,table,xcolor={dvipsnames}]{beamer}
\usetheme{Pittsburgh}
\usecolortheme{seagull}
%\usepackage[utf8]{inputenc}
\usepackage{fontspec}
\usepackage[brazilian]{babel}
\usepackage{amsmath}
\usepackage{listings}
\usepackage{tikz}
\usetikzlibrary{calc,shapes.multipart,chains,arrows, positioning, automata}
\usepackage{multirow}
\usepackage{amsfonts}
\usepackage{amssymb}
\usepackage{lstlinebgrd}
\usepackage{graphicx}
\author{Listas Circulares}
\title{Estruturas de Dados}
%\setbeamercovered{transparent}
\setbeamertemplate{navigation symbols}{}
%\logo{\includegraphics[scale=0.015]{Brasao_UFSC.png}\includegraphics[scale=0.2]{brasao_PPGCC.jpg}}
\institute{Departamento de Informática e de Estatística \\ Prof. Jean Everson Martina \\ Prof. Aldo von Wangenheim}
\date{2016.2}
\subject{}
\usebackgroundtemplate{\includegraphics[width=\paperwidth,
height=\paperheight]{../reusable_images/fundo_UFSC.png}}
\begin{document}

{
\usebackgroundtemplate{\includegraphics[width=\paperwidth,
height=\paperheight]{../reusable_images/fundo_capa.png}}
\begin{frame}
\titlepage
\includegraphics[scale=0.3]{../reusable_images/brasao_INE.png}
\end{frame}
}

\begin{frame}[fragile]{Para que servem listas circulares?}
\lstset{language=C++,
          keywordstyle=\color{blue}\ttfamily,
          stringstyle=\color{red}\ttfamily,
          commentstyle=\color{OliveGreen}\ttfamily,
          breaklines=true,
          basicstyle=\ttfamily\footnotesize
          }
\begin{itemize}
\item Modelagem e planejamento de rotas circulares;
\item Caminho de retorno é diferente do de ida;
\item Exemplo: rotas aéreas tipo volta-ao-mundo;
\item Útil em sistemas de planejamento de itinerário e de reserva global de passagens como Amadeus Global Travel Distribution System
\end{itemize}
\end{frame}

\begin{frame}[fragile]{Para que servem listas circulares?}
\lstset{language=C++,
          keywordstyle=\color{blue}\ttfamily,
          stringstyle=\color{red}\ttfamily,
          commentstyle=\color{OliveGreen}\ttfamily,
          breaklines=true,
          basicstyle=\ttfamily\footnotesize
          }
\begin{itemize}
\item Escalonamento de Processos (Process Scheduling) em kernels de sistemas operacionais multitarefa preemptivos;
\item Exemplo: método round-robin de escalonamento de processos com fatia de tempo constante por processo;\pause
\item Curiosidade: tradução de round-robin é sabiá-cíclico.
\end{itemize}
\end{frame}

\begin{frame}[fragile]{Para que servem listas circulares?}
\lstset{language=C++,
          keywordstyle=\color{blue}\ttfamily,
          stringstyle=\color{red}\ttfamily,
          commentstyle=\color{OliveGreen}\ttfamily,
          breaklines=true,
          basicstyle=\ttfamily\footnotesize
          }
\begin{itemize}
\item Sistemas de impressão que controlam várias impressoras ;
\item Realizam escalonamento de impressoras (round-robin printing queue) por disponibilidade para impressão de documentos em uma fila de impressão cíclica única para várias impressoras.
\end{itemize}
\end{frame}

\begin{frame}[fragile]{Modelagem de Lista Encadeada Circular}
\setbeamercovered{invisible}
\begin{itemize}
\item Último elemento conecta ao primeiro!
\end{itemize}
\begin{center}
\begin{tikzpicture}[list/.style={rectangle split, rectangle split parts=2,    draw=green!60, fill=green!5, rectangle split horizontal}, 
    element/.style={rectangle, draw=yellow!60, fill=yellow!5,},    
    >=stealth, start chain]
  \node[list,on chain] (H) {3};
  \node[list,on chain] (A) [below=of H] {};
  \draw[*->] let \p1 = (H.two), \p2 = (H.center) in (\x1+2,\y2+2) -- (A);
  \node[list,on chain] (B) {};
  \draw[*->] let \p1 = (A.two), \p2 = (A.center) in (\x1,\y2) -- (B);
  \node[list,on chain] (C) {};
  \draw[*->] let \p1 = (B.two), \p2 = (B.center) in (\x1,\y2) -- (C);
  \draw[*->]let \p1 = (C.two), \p2 = (C.center) in (\x1,\y2+2) [out=-50, in=-150, looseness = 3] to (A.west); 
  \node[element] (AP) [below=of A] {melão};
  \draw[*->] let \p1 = (A.one), \p2 = (A.center) in (\x1+2,\y2+2) -- (AP);
  \node[element] (BP) [below=of B] {maçã};
  \draw[*->] let \p1 = (B.one), \p2 = (B.center) in (\x1+2,\y2+2) -- (BP);
  \node[element] (CP) [below=of C] {uva};
  \draw[*->] let \p1 = (C.one), \p2 = (C.center) in (\x1+2,\y2+2) -- (CP);
 \end{tikzpicture}
\end{center}
\end{frame}

\begin{frame}[fragile]{Lista Encadeada Circular Unitária}
\setbeamercovered{invisible}
\begin{center}
\begin{tikzpicture}[list/.style={rectangle split, rectangle split parts=2,
    draw=green!60, fill=green!5, rectangle split horizontal}, 
    element/.style={rectangle, draw=yellow!60, fill=yellow!5,},    
    >=stealth, start chain]

  \node[list,on chain] (H) {3};
  \node[list,on chain] (A) [below=of H] {};
  \draw[*->] let \p1 = (H.two), \p2 = (H.center) in (\x1+2,\y2+2) -- (A);
  \draw[*->]let \p1 = (A.two), \p2 = (A.center) in (\x1,\y2+2) [out=-50, in=-150, looseness = 4] to (A.west); 
  \node[element] (AP) [below=of A] {melão};
  \draw[*->] let \p1 = (A.one), \p2 = (A.center) in (\x1+2,\y2+2) -- (AP);
 \end{tikzpicture}\pause
\end{center}
\begin{itemize}
\item Quais operações têm algoritmos diferentes?\pause
\item Qual é a posição mais indicada para inserir quando a posição não importa ?
\end{itemize}
\end{frame}


\begin{frame}[fragile]{Lista Encadeada Circular - Solução Fácil}
\setbeamercovered{invisible}
\begin{center}
\begin{tikzpicture}[list/.style={rectangle split, rectangle split parts=2,    draw=green!60, fill=green!5, rectangle split horizontal}, 
    element/.style={rectangle, draw=yellow!60, fill=yellow!5,},    
    >=stealth, start chain]
  \node[list,on chain] (H) {3};
  \node[list,on chain,draw=blue!60, fill=blue!5] (D) [below=of H] {};
  \node[list,on chain] (A) {};  
  \draw[*->] let \p1 = (H.two), \p2 = (H.center) in (\x1+2,\y2+2) -- (D);
  \node[list,on chain] (B) {};
  \draw[*->] let \p1 = (A.two), \p2 = (A.center) in (\x1,\y2) -- (B);
  \node[list,on chain] (C) {};
  \draw[*->] let \p1 = (B.two), \p2 = (B.center) in (\x1,\y2) -- (C);

  \draw[*->] let \p1 = (D.two), \p2 = (D.center) in (\x1,\y2) -- (A);
  \node[on chain,draw=purple!60, fill=purple!5,inner sep=6pt] (E) [below =of D] {null*};
  \draw[*->] let \p1 = (D.one), \p2 = (D.center) in (\x1+2,\y2+2) -- (E);
  \draw[*->]let \p1 = (C.two), \p2 = (C.center) in (\x1,\y2+2) [out=-50, in=-150, looseness = 2] to (D.west); 
  \node[element] (AP) [below=of A] {melão};
  \draw[*->] let \p1 = (A.one), \p2 = (A.center) in (\x1+2,\y2+2) -- (AP);
  \node[element] (BP) [below=of B] {maçã};
  \draw[*->] let \p1 = (B.one), \p2 = (B.center) in (\x1+2,\y2+2) -- (BP);
  \node[element] (CP) [below=of C] {uva};
  \draw[*->] let \p1 = (C.one), \p2 = (C.center) in (\x1+2,\y2+2) -- (CP);
 \end{tikzpicture}
\end{center}
\end{frame}


\begin{frame}[fragile]{Lista Encadeada Circular - Solução Fácil}
\setbeamercovered{invisible}
\begin{itemize}
\item Alternativa para evitar mudança nos algoritmos;
\item Nodo-sentinela funciona sempre como nodo sem informações anexado à cabeça de lista;
\item A lista é criada posuindo o nodo sentinela;
\item Os algoritmos não se modificam.
\end{itemize}
\end{frame}

\begin{frame}[fragile]{Trabalho Lista Encadeada Circular}
\lstset{language=C++,
          keywordstyle=\color{blue}\ttfamily,
          stringstyle=\color{red}\ttfamily,
          commentstyle=\color{OliveGreen}\ttfamily,
          breaklines=true,
          basicstyle=\ttfamily\footnotesize
          }
\begin{itemize}
\item Implemente uma classe Lista todas as operações vistas;
\item Implemente a lista usando Templates;
\item Use as melhores práticas de orientação a objetos;
\item Documente todas as classes, métodos e atributos;
\item Aplique os testes unitários disponíveis no moodle da disciplina para validar sua estrutura de dados;
\item Entregue até a data definida no moodle.
\end{itemize}
\end{frame}

\begin{frame}[fragile]{Lista Duplamente Encadeada Circular}
\setbeamercovered{invisible}
\begin{itemize}
\item É muito simples. Ao invés de apontar o inicio e o fim pra null*, apontamos um para o outro;
\item Não opecisa de sentinela;
\item Basta uma modificação no adicionaNoInicio e no retiraDoInicio;
\end{itemize}
\end{frame}

\begin{frame}[fragile]{Trabalho Lista Duplamente Encadeada Circular}
\lstset{language=C++,
          keywordstyle=\color{blue}\ttfamily,
          stringstyle=\color{red}\ttfamily,
          commentstyle=\color{OliveGreen}\ttfamily,
          breaklines=true,
          basicstyle=\ttfamily\footnotesize
          }
\begin{itemize}
\item Implemente uma classe Lista todas as operações vistas;
\item Implemente a lista usando Templates;
\item Use as melhores práticas de orientação a objetos;
\item Documente todas as classes, métodos e atributos;
\item Aplique os testes unitários disponíveis no moodle da disciplina para validar sua estrutura de dados;
\item Entregue até a data definida no moodle.
\end{itemize}
\end{frame}

{
\usebackgroundtemplate{\includegraphics[width=\paperwidth,
height=\paperheight]{../reusable_images/fundo_capa.png}}
\begin{frame}

{\LARGE Perguntas????}

\end{frame}
}


{
\usebackgroundtemplate{\includegraphics[width=\paperwidth,
height=\paperheight]{../reusable_images/fundo_capa.png}}
\begin{frame}
\includegraphics[scale=0.8]{../reusable_images/cc_logo_arge.png}\hspace{0.5cm}
\includegraphics[scale=0.95]{../reusable_images/by.png}

\vspace{1cm}
Este trabalho está licenciado sob uma Licença Creative Commons Atribuição 4.0 Internacional. Para ver uma cópia desta licença, visite http://creativecommons.org/licenses/by/4.0/.

\end{frame}
}
\end{document}
